\section{Experimental Setup}

This section details the experimental configuration for validating Malak Platform across three reference applications.

\subsection{Application 1: Vision Jellyfish}

\subsubsection{Task Description}
Real-time detection and classification of marine life (specifically jellyfish species) from underwater camera footage. This application targets autonomous underwater vehicles (AUVs) and stationary monitoring systems in marine biology research.

\subsubsection{Dataset}
\begin{itemize}
\item \textbf{Source}: Custom dataset collected from Mediterranean Sea expeditions
\item \textbf{Size}: 15,000 annotated images (12,000 train, 1,500 validation, 1,500 test)
\item \textbf{Classes}: 8 jellyfish species + background
\item \textbf{Resolution}: 320$\times$240 RGB images
\item \textbf{Challenges}: Variable lighting, occlusion, turbidity, class imbalance
\end{itemize}

\subsubsection{Model Architecture}
\begin{itemize}
\item \textbf{Baseline}: MobileNetV2-SSDLite (3.2M parameters)
\item \textbf{Student}: Custom lightweight detector (680K parameters)
\item \textbf{Backbone}: Depth-multiplier 0.35 MobileNetV2
\item \textbf{Detection head}: Single-scale SSD with 3 anchor boxes
\end{itemize}

\subsubsection{Training Configuration}
\begin{verbatim}
Optimizer: AdamW (lr=1e-3, weight_decay=1e-4)
Batch size: 32
Epochs: 150 with cosine annealing
Augmentation: Random flip, crop, color jitter,
              Gaussian blur, brightness adjustment
Knowledge distillation: Temperature=4.0,
                        alpha=0.7 (soft labels)
\end{verbatim}

\subsubsection{Deployment Target}
STM32H7 microcontroller (Cortex-M7 @ 480MHz, 1MB RAM) with camera interface operating in battery-powered enclosures.

\subsection{Application 2: Energy Optimization}

\subsubsection{Task Description}
Short-term energy consumption forecasting for smart building management systems. Predicts next-hour electricity demand based on historical usage, time-of-day, occupancy sensors, and weather data.

\subsubsection{Dataset}
\begin{itemize}
\item \textbf{Source}: UCI Appliances Energy Prediction dataset
\item \textbf{Size}: 19,735 samples (4.5 months of 10-minute readings)
\item \textbf{Features}: 29 inputs (temperature, humidity, lighting, appliance power)
\item \textbf{Target}: Total energy consumption (kWh)
\item \textbf{Split}: 70\% train, 15\% validation, 15\% test (temporal split)
\end{itemize}

\subsubsection{Model Architecture}
\begin{itemize}
\item \textbf{Baseline}: LSTM with 2 layers (128 units each, 265K parameters)
\item \textbf{Student}: Quantized LSTM with 1 layer (64 units, 41K parameters)
\item \textbf{Sequence length}: 24 timesteps (4 hours history)
\item \textbf{Output}: Single-step ahead prediction
\end{itemize}

\subsubsection{Training Configuration}
\begin{verbatim}
Optimizer: Adam (lr=5e-4)
Batch size: 64
Epochs: 100 with early stopping (patience=15)
Loss: Huber loss (delta=1.0)
Normalization: Min-max scaling per feature
Knowledge distillation: MSE between teacher and
                        student hidden states
\end{verbatim}

\subsubsection{Deployment Target}
Raspberry Pi 4 (Cortex-A72 @ 1.5GHz) integrated with building automation system, updating predictions every 10 minutes.

\subsection{Application 3: Neuro MS}

\subsubsection{Task Description}
Segmentation of multiple sclerosis (MS) lesions in brain MRI scans for clinical decision support. Privacy-critical application requiring on-device processing to comply with HIPAA/GDPR.

\subsubsection{Dataset}
\begin{itemize}
\item \textbf{Source}: ISBI 2015 Longitudinal MS Lesion Segmentation Challenge
\item \textbf{Size}: 14 patients, 5 timepoints each (70 scans total)
\item \textbf{Modalities}: T1-weighted, T2-weighted, FLAIR
\item \textbf{Resolution}: 181$\times$217$\times$181 voxels
\item \textbf{Preprocessing}: Skull stripping, bias field correction, registration
\item \textbf{Augmentation}: Random rotation, elastic deformation, intensity shifts
\end{itemize}

\subsubsection{Model Architecture}
\begin{itemize}
\item \textbf{Baseline}: 3D U-Net (8.6M parameters)
\item \textbf{Student}: Efficient U-Net with depthwise separable convolutions (1.2M parameters)
\item \textbf{Depth}: 4 encoder/decoder levels
\item \textbf{Filters}: [16, 32, 64, 128] with bottleneck at 256
\end{itemize}

\subsubsection{Training Configuration}
\begin{verbatim}
Optimizer: SGD with Nesterov momentum (lr=1e-2)
Batch size: 2 (3D patches of 64^3)
Epochs: 200
Loss: Combined Dice + cross-entropy
Class weighting: 10:1 (lesion:background)
Knowledge distillation: Feature matching at
                        decoder layers
\end{verbatim}

\subsubsection{Deployment Target}
NVIDIA Jetson Nano (128-core GPU) for edge inference in clinical workstations, processing full 3D volumes in <5 seconds.

\subsection{Compression Experiments}

For each application, we evaluate:

\subsubsection{Quantization Sweep}
\begin{itemize}
\item FP32 (baseline)
\item INT8 post-training quantization (PTQ)
\item INT8 quantization-aware training (QAT)
\item INT4 weights + INT8 activations
\item Mixed precision (automatic bit-width search)
\end{itemize}

\subsubsection{Pruning Configurations}
\begin{itemize}
\item Magnitude-based channel pruning at 30\%, 50\%, 70\% sparsity
\item Iterative pruning with fine-tuning (5 cycles)
\item Combined quantization + pruning
\end{itemize}

\subsubsection{Compiler Backends}
\begin{itemize}
\item TVM with AutoTVM tuning (1000 trials per operator)
\item MLIR with hand-tuned quantization passes
\item XLA (for Coral Edge TPU only)
\item Baseline C++ (no compiler optimization)
\end{itemize}

\subsection{Profiling Setup}

\subsubsection{Latency Measurement}
\begin{itemize}
\item 1000 inferences with warm-up (50 iterations)
\item Per-layer profiling using on-chip cycle counters
\item Reported metrics: mean, median, 95th percentile
\end{itemize}

\subsubsection{Energy Measurement}
\begin{itemize}
\item External power monitor (Monsoon Power Monitor) at 5 kHz sampling
\item Isolated inference execution (no concurrent tasks)
\item Energy = $\int$ Power $\times$ dt over inference window
\end{itemize}

\subsubsection{Memory Profiling}
\begin{itemize}
\item Peak RAM via linker map analysis + runtime allocator hooks
\item Flash size from compiled binary (text + data sections)
\item Memory bandwidth via cache miss counters (where available)
\end{itemize}

\subsection{Baseline Comparisons}

We compare Malak Platform against:
\begin{itemize}
\item \textbf{TensorFlow Lite Micro}: Official TFLite interpreter for microcontrollers
\item \textbf{CMSIS-NN}: ARM-optimized kernels with manual graph execution
\item \textbf{MicroTVM}: Apache TVM with default autotuning
\item \textbf{PyTorch Mobile}: PyTorch runtime for mobile devices
\end{itemize}

For fair comparison:
\begin{enumerate}
\item All frameworks use identical trained models (converted to respective formats)
\item Same quantization scheme (INT8 symmetric per-tensor)
\item Same target hardware and compiler toolchains (GCC 10.2, -O3 optimization)
\item No framework-specific optimizations (e.g., operator fusion disabled for baselines lacking it)
\end{enumerate}

\subsection{Statistical Validation}

All performance measurements report:
\begin{itemize}
\item Mean $\pm$ standard deviation across 5 independent runs
\item Statistical significance via paired t-tests (p < 0.05)
\item Confidence intervals (95\%) for accuracy metrics
\end{itemize}
